% chapter/06_experimental_design.tex
% Session 6 — Experimental Design
% Compiled with pdflatex + bibtex (nonstopmode)

\section{Experimental Design}
\label{sec:exp_design}

This section describes the datasets, baselines, evaluation metrics, and implementation
details used throughout the experiments. All results are reported in
Section~\ref{sec:results}.

\subsection{Research Questions}
\label{sec:rq}

The experiments address five questions:

\begin{enumerate}[label=(\textbf{RQ\arabic*)}]
  \item \textbf{SIR sensitivity.} How does DCA-Trie's oracle configuration affect the
        Semantic Irrelevance Ratio across different KG subgraphs?
  \item \textbf{Answer quality.} Does tighter structural constraint from the TypeOracle
        improve Hits@1, or does it leave answer accuracy unchanged?
  \item \textbf{Decomposed contribution.} How much of the SIR reduction comes from the
        type gate versus the range gate?
  \item \textbf{Latency.} What is the overhead of step-wise oracle evaluation relative to
        GCR's static trie construction?
  \item \textbf{v1 vs v2.} When does step-wise dynamic expansion (v2) outperform static
        filtering (v1)?
\end{enumerate}

RQ1 and RQ2 address the core question: whether constraint tightness and answer
correctness are coupled or independent. RQ3 identifies which component drives the
reduction. RQ4 and RQ5 address practical deployment trade-offs.

\subsection{Datasets}
\label{sec:datasets}

We evaluate on WebQuestionSP (WebQSP) \citep{yih2016websqp} and Complex WebQuestions
(CWQ) \citep{talmor2018cwq}. Both benchmark datasets query Freebase
\citep{bollacker2008freebase}, a large-scale general-domain knowledge graph.

WebQSP contains 5,813 questions with an average of 2.1 entities per answer and 2 hops
of reasoning depth. We use the standard train/test split (3,034 train, 2,502 test; we
report on the test split with 1,628 questions following GCR's partition). CWQ contains
34,211 questions requiring multi-hop reasoning over Freebase, with more complex
compositional structure. We use the standard split (27,639 train, 3,531 test).

Freebase contains 122 million entities, 10,904 relations, and 1.5 billion triples. We
use the same Freebase snapshot as GCR \citep{luo2025-graph-constrained-reasoning}. For
each question, we extract all entities mentioned in the question, enumerate all 2-hop
paths from those entities, and index them in a trie for constrained decoding.

\subsection{Baselines}
\label{sec:baselines}

We compare DCA-Trie against two baseline categories.

\textbf{LLM reasoning.} Unconstrained Llama-3.1-8B \citep{touvron2023llama} generates
answers without KG grounding. This establishes a lower bound: what the model achieves
without any structural constraint.

\textbf{GCR} \citep{luo2025-graph-constrained-reasoning}. The direct baseline. GCR
constructs a static trie from all 2-hop paths before decoding begins and constrains
beam search to paths present in the trie. Our re-implementation achieves 91.6\% Hits@1
on WebQSP, matching the reported result. We use the same trie construction, beam
parameters, and LLM checkpoint.

\textbf{Why these two?} The unconstrained baseline tells us whether KG grounding helps
at all. GCR tells us whether dynamic constraint adjustment adds value beyond static
constraint. Comparing against GCR is the primary evaluation: DCA-Trie modifies only
the constraint oracle, keeping everything else identical.

\subsection{Evaluation Metrics}
\label{sec:metrics}

\textbf{Hits@1.} Whether the top-1 generated answer appears in the set of gold answers.
This is the standard metric for KGQA benchmarks \citep{luo2025-graph-constrained-reasoning,
li2024-decoding-on-graphs}. A higher value indicates better answer accuracy.

\textbf{SIR (Semantic Irrelevance Ratio).} The fraction of structurally valid trie
paths that are semantically irrelevant to the question (Definition~\ref{def:sir}). A
lower value indicates tighter constraint. SIR decouples constraint quality from answer
accuracy: two systems can have identical Hits@1 but different SIR values.

\textbf{FNR (False Negative Rate).} The fraction of gold-answer paths that the oracle
incorrectly removes. An oracle with FNR$>0$ excludes correct answers before the LLM
even sees them. We report FNR for the type gate and range gate separately.

\textbf{Latency.} Wall-clock time per question, measured from the start of decoding to
the generation of the final token. We report the mean and 95\% confidence interval over
the test split.

\textbf{Why these metrics?} Hits@1 measures what the system gets right. SIR measures
what the oracle filters. FNR measures what the oracle removes that it should not have.
Latency measures the cost of the dynamic oracle. Together they answer whether tighter
constraints actually help, and at what cost.

\subsection{Implementation Details}
\label{sec:implementation}

\textbf{Model.} We use \texttt{rmanluo/GCR-Meta-Llama-3.1-8B-Instruct}, a Llama-3.1-8B
model fine-tuned on GCR's KG-reasoning training data. This is the same checkpoint used
in GCR's reported results.

\textbf{Hardware.} All experiments run on a single NVIDIA RTX 4090 (24\,GB VRAM) with
bfloat16 precision.

\textbf{Decoding.} We use group beam search with $k=10$ beams, a diversity penalty of
1.0, and a maximum generation length of 256 tokens. The index path length is 2 hops,
matching GCR's configuration.

\textbf{TypeOracle.} The type gate uses a regex-based question classifier that maps
question patterns (e.g., ``what is the capital of X'' $\rightarrow$ \texttt{Location})
to Freebase types. The range gate uses relation metadata from Freebase's schema. Both
gates require no training and no neural inference.

\textbf{Reproducibility.} The TypeOracle stores ontology metadata ($<2$\,MB for
Freebase). The regex classifier is deterministic. All random seeds are fixed. The
constrained decoding pipeline is identical to GCR except for the oracle, ensuring that
any performance differences are attributable to the oracle alone.